\documentclass[10pt,a4paper]{amsart}
\usepackage[margin=19mm]{geometry}
\usepackage{amsmath,amssymb,mathtools}
\usepackage[scr=rsfs,cal=euler]{mathalfa}
\usepackage{xfrac}
\usepackage[unicode,hidelinks,bookmarks=false]{hyperref}
\newcommand{\ie}{i.e.\ }
\newcommand{\cf}{cf.\ }
\newcommand{\resp}{respectively\ }
\newcommand{\Subsection}{\subsection}
\providecommand{\nobreakdash}{\nobreak\hbox{-}}
\setlength{\emergencystretch}{2em}
\multlinegap0pt
\usepackage{bm}
\usepackage{bbm}
\usepackage{dsfont}
\usepackage{xspace}
\usepackage{cancel}
\usepackage{mathtools}
\usepackage{amsthm}
\makeatletter
\@ifundefined{theorem}{\newtheorem{theorem}{Theorem}}{}
\@ifundefined{lemma}{\newtheorem{lemma}[theorem]{Lemma}}{}
\@ifundefined{proposition}{\newtheorem{proposition}[theorem]{Proposition}}{}
\makeatother
\newcommand{\R}{\mathbb{R}}
\newcommand{\cG}{\mathcal{G}}
\title[On the Schiffer and Berenstein conjectures with high-frequen]{On the Schiffer and Berenstein conjectures with high-frequency for convex domains in the plane}
\author{Guowei Dai, Yingxin Sun, Juncheng Wei, Yong Zhang}
\date{Source: arXiv:2511.19819v2 (CC BY 4.0)}
\begin{document}
\maketitle

\noindent\emph{Source and licence.} On the Schiffer and Berenstein conjectures with high-frequency for convex domains in the plane. Version arXiv:2511.19819v2 (CC BY 4.0), distributed under the Creative Commons Attribution 4.0 International licence, as recorded in the arXiv licence metadata for this exact version and shown on the abstract page https://arxiv.org/abs/2511.19819v2. Full rights notice: licenses/arxiv-2511.19819-CC-BY-4.0.txt.\par\smallskip

\noindent\textit{Editorial scope.} Counted unit: the Proposition whose source label is \texttt{prop:defect-distance} in the article (source0.tex lines 661--672, source label \texttt{prop:defect-distance}), delivered together with the article's own complete original proof of it (source0.tex lines 674--801, one \texttt{proof} environment of 128 lines). No mathematics is written, repaired or re-derived: statement and proof are the article's own text line for line. Required context retained uncounted: eq:rho-bounds (lines 524--526); lem:periodic-interpolation (lines 634--644). Every definition, display and label of those lines is retained unchanged. Omitted from this package and not redistributed: 1--523, 527--633, 645--660, 673--673, 802--1797. Nothing inside a retained excerpt is omitted or altered: every retained body is delivered exactly as the article's lines are written, so no omission receipt of any kind accompanies this package. Citations: the retained excerpts carry no citation command at all, so no bibliography entry of the article is transcribed and no reference is redistributed. Reference adaptation: none was needed and none was made; every reference inside the delivered unit resolves inside it, so no unresolved reference or citation is produced. Printed slips kept literal: no printed word, symbol, hypothesis or number of the counted statement or of its proof was altered. Layout and class: the article's own class is \texttt{article} with packages including mathrsfs, amssymb, amsmath, amsfonts, amstext, amsthm, graphicx, subfig, color, indentfirst, latexsym, bm, cite, xy; the wrapper is the lane's pinned \texttt{amsart} editorial wrapper at 10pt on A4, which restates the theorem-like environments the retained text needs and adds the article's own macro definitions that the retained text needs (\texttt{\textbackslash R}, \texttt{\textbackslash cG}), with extra wrapper packages (bm, bbm, dsfont, xspace, cancel, mathtools). The delivered numbering is the unit's own. Assets: no retained span contains a figure environment, a graphics-inclusion command, a PDF-page inclusion, a table or a TikZ picture; no image file is copied into this package, the package declares no asset dependency and no image file is redistributed with it. Licence: version arXiv:2511.19819v2 of this article is distributed under the Creative Commons Attribution 4.0 International licence, as recorded in the arXiv licence metadata for this exact version and shown on the abstract page https://arxiv.org/abs/2511.19819v2. A full rights notice is delivered at licenses/arxiv-2511.19819-CC-BY-4.0.txt. Characters: every retained excerpt is pure ASCII, so the delivered engine, XeLaTeX with its default Latin Modern text font, reports no missing character.
\begin{equation}\label{eq:rho-bounds}
\frac1{\kappa_1}\le\rho(\theta)\le\frac1{\kappa_0}.
\end{equation}

\begin{lemma}[Periodic interpolation]\label{lem:periodic-interpolation}
For every integer $m\geq3$ there exists $C_m>0$ such that every
$u\in C^m\left(\R/2\pi\mathbb Z\right)$ satisfies
\begin{equation}\label{eq:periodic-interpolation}
 \left\|u''\right\|_{L^\infty}
 \leq C_m\left(
 \left\|u'\right\|_{L^\infty}^{(m-2)/(m-1)}
 \left\|u^{(m)}\right\|_{L^\infty}^{1/(m-1)}
 +\left\|u'\right\|_{L^\infty}\right).
\end{equation}
\end{lemma}

\begin{proposition}[Regularity-dependent defect estimate]
\label{prop:defect-distance}
Let $m\geq5$ be an integer. For every $M,\kappa_0,\kappa_1>0$ there exists
$c_m=c_m\left(M,\kappa_0,\kappa_1\right)>0$ such that, for every $R>0$ and every
$\Omega\in\cG_m\left(R,M,\kappa_0,\kappa_1\right)$,
\begin{equation}\label{eq:defect-distance-lower}
 \mathfrak B(\Omega)
 \geq c_m
 \min\left\{d_{\mathrm{disk}}(\Omega)^{p_m},1\right\}.
\end{equation}
The constant is independent of $R$ and of the particular domain.
\end{proposition}

\begin{proof}
Write $\beta=\mathfrak B(\Omega)$ and abbreviate
$\rho_\pi(\theta)=\rho(\theta+\pi)$ and
$h'_\pi(\theta)=h'(\theta+\pi)$. The four quantities
\[
 X_\pm=\frac12\log\rho(\theta)\pm h'(\theta),
 \qquad
 Y_\pm=\frac12\log\rho_\pi(\theta)\mp h'_\pi(\theta)
\]
range in a compact interval determined only by
$M,\kappa_0,$ and $\kappa_1$. By the definition of $\beta$,
\begin{equation}\label{zhishubound}
 \left|e^{X_\pm}-e^{Y_\pm}\right|\leq\beta.
\end{equation}
Here and below, the constants $C_j$ depend only on
$m,M,\kappa_0,$ and $\kappa_1$.
All four exponentials are bounded below by
$\mu=\kappa_1^{-1/2}e^{-M}>0$.
The logarithmic function $x \mapsto \log x$ is Lipschitz on $[\mu, \infty)$ with constant $\mu^{-1}$.
Therefore, from (\ref{zhishubound}),
\begin{align*}
\left|X_+ - Y_+\right|
&= \left| \log e^{X_+} - \log e^{Y_+} \right|
\le \frac{1}{\mu} \left|e^{X_+} - e^{Y_+}\right|
\le \frac{\beta}{\mu}, \\
\left|X_- - Y_-\right|
&= \left| \log e^{X_-} - \log e^{Y_-} \right|
\le \frac{1}{\mu} \left|e^{X_-} - e^{Y_-}\right|
\le \frac{\beta}{\mu}.
\end{align*}
Adding and subtracting the two resulting inequalities give
\[
\left| \log \rho(\theta) - \log \rho(\theta+\pi) \right| \le \frac{2\beta}{\mu}.
\]
and
\[
\left| h'(\theta) + h'(\theta+\pi) \right| \le \frac{\beta}{\mu}.
\]
Since $\theta$ is arbitrary, taking the $L^\infty$ norm over $\theta \in \mathbb{R}/2\pi\mathbb{Z}$ gives
\[
\left\| h'(\cdot) + h'(\cdot+\pi)\right\|_{L^\infty} \le \frac{1}{\mu}\beta
\]
and
\[
\left\| \log \rho(\cdot) - \log \rho(\cdot+\pi)\right\|_{L^\infty} \le \frac{2}{\mu}\beta.
\]
Adding the two inequalities yields
\begin{equation}\label{eq:log-defect-control}
 \left\|h'(\cdot)+h'(\cdot+\pi)\right\|_{L^\infty}
 +\left\|\log\rho(\cdot)-\log\rho(\cdot+\pi)\right\|_{L^\infty}
 \leq C_0\beta.
\end{equation}

From the curvature bounds in \eqref{eq:rho-bounds}, we have that
both $\rho(\theta)$ and $\rho(\theta+\pi)$ lie in the compact interval
$\left[1/\kappa_1, 1/\kappa_0\right]$ for any fixed $\theta$.
By the mean value theorem applied to $\log x$ on the interval $\left[1/\kappa_1, 1/\kappa_0\right]$,
there exists $\xi$ between $\rho(\theta)$ and $\rho(\theta+\pi)$ such that
\[
\log \rho(\theta) - \log \rho(\theta+\pi) = \frac{1}{\xi}(\rho(\theta)-\rho(\theta+\pi)).
\]
Taking absolute values gives
\[
\left|\rho(\theta)-\rho(\theta+\pi)\right| = \xi |\log \rho(\theta) - \log \rho(\theta+\pi)|.
\]
From (\ref{eq:log-defect-control}) and the fact that $\xi \le 1/\kappa_0$,
we have
\begin{equation*}
|\rho(\theta)-\rho(\theta+\pi)| \le \frac{1}{\kappa_0} |\log \rho(\theta) - \log\rho(\theta+\pi)|\le \frac{C_0}{\kappa_0} \beta.
\end{equation*}
Since the above inequality holds for every $\theta$, taking the supremum over
$\theta \in \mathbb{R}/2\pi\mathbb{Z}$ yields
\begin{equation}\label{eq:rho-defect-control}
 \|\rho(\cdot)-\rho(\cdot+\pi)\|_{L^\infty}
 \leq C_1\beta.
\end{equation}

Decompose $h$ into its antipodally even and odd parts,
\[
 h_{\mathrm{e}}(\theta)=\frac{h(\theta)+h(\theta+\pi)}{2},
 \qquad
 h_{\mathrm{o}}(\theta)=\frac{h(\theta)-h(\theta+\pi)}{2}.
\]
Equations \eqref{eq:log-defect-control}--\eqref{eq:rho-defect-control}
give
\begin{equation}\label{eq:even-odd-control}
 \left\|h_{\mathrm{e}}'\right\|_{L^\infty}\leq C_2\beta,
 \qquad
 \left\|h_{\mathrm{o}}+h_{\mathrm{o}}''\right\|_{L^\infty}\leq C_2\beta.
\end{equation}
Steiner normalization means that the first sine and cosine Fourier
coefficients of $h$, and hence of $h_{\mathrm{o}}$, vanish. The kernel of
$1+\partial_\theta^2$ on $2\pi$-periodic functions is precisely the span of
$\cos\theta$ and $\sin\theta$. The elementary periodic Green-function
estimate on the complementary subspace therefore yields
\begin{equation}\label{eq:odd-C2-control}
 \left\|h_{\mathrm{o}}\right\|_{C^2}\leq
 C_3\left\|h_{\mathrm{o}}+h_{\mathrm{o}}''\right\|_{L^\infty}
 \leq C_4\beta.
\end{equation}

Let $r=(2\pi)^{-1}\int_0^{2\pi}h_{\mathrm{e}}(\theta)\,\mathrm{d}\theta$; this number is
positive because the Steiner point lies in the interior of $\Omega$.
The periodic Poincar\'e inequality and \eqref{eq:even-odd-control} give
\begin{equation*}\label{eq:even-C1-control}
 \left\|h_{\mathrm{e}}-r\right\|_{C^1}\leq C_5\beta.
\end{equation*}
Applying Lemma \ref{lem:periodic-interpolation} to $h_{\mathrm{e}}$ gives
\begin{align}
 \left\|h_{\mathrm{e}}''\right\|_{L^\infty}
 &\leq C_6\left(
 \left\|h_{\mathrm{e}}'\right\|_{L^\infty}^{(m-2)/(m-1)}
 \left\|h_{\mathrm{e}}^{(m)}\right\|_{L^\infty}^{1/(m-1)}
 +\left\|h_{\mathrm{e}}'\right\|_{L^\infty}\right)\nonumber\\
 &\leq C_7\left(\beta^{1/p_m}+\beta\right).
 \label{eq:even-C2-interpolation}
\end{align}
Combining \eqref{eq:odd-C2-control}--\eqref{eq:even-C2-interpolation} and $h=h_{\mathrm{e}}+h_{\mathrm{o}}$ gives
\begin{equation*}\label{eq:distance-upper-by-defect}
 d_{\mathrm{disk}}(\Omega)
 \leq\|h-r\|_{C^2}
 \leq C_8\left(\beta^{1/p_m}+\beta\right).
\end{equation*}
If $\beta\leq1$, this implies
$\beta\geq\left(2C_8\right)^{-p_m}d_{\mathrm{disk}}(\Omega)^{p_m}$; if
$\beta>1$, the required estimate is immediate after decreasing the
constant. This proves \eqref{eq:defect-distance-lower}.
\end{proof}

\end{document}
